% F5b -- Full fine-tuning, LoRA and QLoRA, drawn side by side.
\begin{figure}[!htbp]
\centering
\fitwidth{%
\begin{tikzpicture}[font=\small,
  big/.style={draw=clrule,fill=clnonsig!30,minimum width=2.5cm,minimum height=2.5cm},
  small/.style={draw=clgain,fill=clgain!25,minimum width=2.5cm,minimum height=0.42cm},
  smallv/.style={draw=clgain,fill=clgain!25,minimum width=0.42cm,minimum height=2.5cm},
  q/.style={draw=clmodel,fill=clmodel!18,minimum width=2.5cm,minimum height=2.5cm},
  lbl/.style={font=\scriptsize,text=clrule,align=center},
]
% ---------------- (a) full fine-tuning ----------------
\node[big,fill=clloss!25,draw=clloss] (fa) at (0,0) {};
\node[lbl,above=3pt of fa] {\textbf{(a) Full fine-tuning}};
\node[lbl,below=3pt of fa,align=center] {every weight is updated\\ \textcolor{clloss}{100\,\% trainable}\\ does not fit on a free GPU};
\node[font=\scriptsize,text=clloss] at (fa) {$W$};

% ---------------- (b) LoRA ----------------
\node[big] (fb) at (4.6,0) {};
\node[font=\scriptsize,text=clrule] at (fb) {$W$ \emph{frozen}};
\node[smallv] (bA) at (6.35,0.55) {};
\node[small]  (bB) at (6.35,-1.0) {};
\node[font=\scriptsize] at (bA) {$A$};
\node[font=\scriptsize] at (bB) {$B$};
\node[lbl,above=3pt of fb] {\textbf{(b) LoRA}};
\node[lbl,below=17pt of fb,align=center,xshift=0.9cm]
  {the frozen matrix is left alone; only the two\\ thin \textcolor{clgain}{$A$} and
   \textcolor{clgain}{$B$} matrices are trained\\ \textcolor{clgain}{$\approx$2\,\% trainable}};

% ---------------- (c) QLoRA ----------------
\node[q] (fc) at (10.4,0) {};
\node[font=\scriptsize,text=clmodel,align=center] at (fc) {$W$ \emph{frozen}\\ \emph{and 4-bit}};
\node[smallv] (cA) at (12.15,0.55) {};
\node[small]  (cB) at (12.15,-1.0) {};
\node[font=\scriptsize] at (cA) {$A$};
\node[font=\scriptsize] at (cB) {$B$};
\node[lbl,above=3pt of fc] {\textbf{(c) QLoRA}, used here};
\node[lbl,below=17pt of fc,align=center,xshift=0.9cm]
  {same, plus the frozen matrix is stored\\ at 4 bits per number instead of 16\\
   \textcolor{clmodel}{$\approx$1/4 of the frozen memory}};
\end{tikzpicture}}

\caption{What LoRA and QLoRA do}
\label{fig:lora}
\figtext{%
\figpara{The grey square is the model.}{A language model is essentially a stack
of very large tables of numbers, and adapting one to a new task means changing
numbers in those tables. The three panels are the three ways of doing it, drawn
at the same scale.}

\figpara{(a) Change everything.}{Every number in the table is modified. This has
the highest ceiling in principle, but while the run lasts each number has to be
held in the graphics card's memory together with about three more used to keep
track of how to adjust it. The whole thing is several times larger than the
model, and it does not fit on a free card.}

\figpara{(b) LoRA: leave the table alone and add two thin ones beside
it.}{\cite{hu2022lora} The big table is frozen and never touched. Two narrow
tables are added next to it, and only those are trained. Two narrow tables
standing in for a change to a wide one is the whole of the idea, and here it
means about 2 numbers in every 100 are ever allowed to move.}

\figpara{(c) QLoRA: the same, with the frozen table compressed.}{\cite{dettmers2023qlora}
This is the method used throughout the report. On top of LoRA, the frozen table
is stored with four bits per number instead of sixteen, which quarters the memory
it occupies. The compression is safe precisely because those numbers are frozen:
they are read, never updated. The added tables stay at full precision.}

\figpara{Why this decides the shape of the whole internship.}{Not accuracy, but
waiting time. With (c), one training run fits in half an hour on a free graphics
card, so nine weeks buy thirty experiments instead of two.}}

\end{figure}
